\section*{Hoofdstuk \ref{ch:g11:scal} --- Het scalair product in het vlak}

\begin{solution}{exo:g11:scal:1}
$(2,5)\cdot(3,-1) = 6 - 5 = 1$.

$(1,-4)\cdot(8,2) = 8 - 8 = 0$ (\omterm{def:g11:scal:orthogonal}{orthogonale vectoren}).

$\vec u\cdot\vec v = 3 \times 4 \times \cos\frac\pi3 = 12 \times \frac12
= 6$.
\end{solution}

\begin{solution}{exo:g11:scal:2}
$\vect{AB}\cdot\vect{AC} = 6 \times 6 \times \cos 60^\circ = 18$.

De \omterm{def:g11:vect:vector}{vectoren} $\vect{AB}$ en $\vect{BC}$ maken een hoek van
$180^\circ - 60^\circ = 120^\circ$ (zet ze staart aan staart in $B$), dus is
$\vect{AB}\cdot\vect{BC} = 36 \cos 120^\circ = -18$.
\end{solution}

\begin{solution}{exo:g11:scal:3}
$\vec u \cdot \vec v = 3(t - 8) + t = 4t - 24 = 0$ geeft $t = 6$.

$\vec u - \vec v = (t - 4,\ 0)$, dus
$\vec u \cdot (\vec u - \vec v) = t(t-4) + 0 = t^2 - 4t$, wat verdwijnt voor
$t = 0$ of $t = 4$.
\end{solution}

\begin{solution}{exo:g11:scal:4}
$\vec u \cdot \vec v = 3 + 2 = 5$, $\norm{\vec u} = \sqrt5$ en
$\norm{\vec v} = \sqrt{10}$, dus
\[
\cos\theta = \frac{5}{\sqrt5\,\sqrt{10}} = \frac{5}{5\sqrt2}
= \frac{\sqrt2}{2},
\]
en bijgevolg is $\theta = 45^\circ$ exact.
\end{solution}

\begin{solution}{exo:g11:scal:5}
$a^2 = 16 + 36 - 2 \times 4 \times 6 \times \cos\frac{2\pi}{3}
= 52 - 48 \times \left(-\frac12\right) = 76$, dus
$a = \sqrt{76} = 2\sqrt{19} \approx 8.7$.

Voor de tweede driehoek geeft de cosinusregel, opgelost naar de hoek,
\[
\cos\widehat A = \frac{b^2 + c^2 - a^2}{2bc}
= \frac{25 + 16 - 49}{2 \times 5 \times 4} = -\frac{8}{40} = -\frac15,
\]
negatief: $\widehat A$ is stomp.
\end{solution}

\begin{solution}{exo:g11:scal:6}
$\vect{AB}\,(4, -2)$ en $\vect{AC}\,(2, 4)$:
$\vect{AB}\cdot\vect{AC} = 8 - 8 = 0$, dus is de hoek in $A$ recht. Zijden:
$AB = \sqrt{16 + 4} = \sqrt{20}$, $AC = \sqrt{4 + 16} = \sqrt{20}$, en
$\vect{BC}\,(-2, 6)$ geeft $BC = \sqrt{40}$. Consistentie:
$BC^2 = 40 = 20 + 20 = AB^2 + AC^2$, de cosinusregel met
$\cos\widehat A = 0$ (Pythagoras).
\end{solution}

\begin{solution}{exo:g11:scal:7}
In middelpunt-straalvorm: $(x + 2)^2 + (y - 3)^2 = 25$.

Voor de tweede cirkel vul je de kwadraten aan:
\[
x^2 - 6x + y^2 + 2y = 6
\iff (x - 3)^2 - 9 + (y + 1)^2 - 1 = 6
\iff (x-3)^2 + (y+1)^2 = 16 :
\]
middelpunt $(3, -1)$, straal $4$.
\end{solution}

\begin{solution}{exo:g11:scal:8}
$M(x,y)$ ligt op de cirkel wanneer $\vect{MA}\cdot\vect{MB} = 0$ met
$\vect{MA}\,(-3 - x,\ -y)$ en $\vect{MB}\,(-x,\ 2 - y)$:
\[
(-3 - x)(-x) + (-y)(2 - y) = x^2 + 3x + y^2 - 2y = 0 .
\]
Voor de oorsprong: $0 + 0 + 0 - 0 = 0$, dus ligt $O$ op de cirkel ---
\omterm{def:g11:seq:geometric}{meetkundig} zijn $\vect{OA}\,(-3,0)$ en $\vect{OB}\,(0,2)$ \omterm{def:g11:scal:orthogonal}{orthogonaal},
zodat $O$ het lijnstuk $[AB]$ onder een rechte hoek ziet. (Kwadraten
aanvullen geeft middelpunt $\left(-\frac32, 1\right)$ en straal
$\frac{\sqrt{13}}{2}$.)
\end{solution}

\begin{solution}{exo:g11:scal:9}
De loodlijn door $P$ heeft de \omterm{def:g11:vect:direction}{richtingsvector} van $d$, namelijk $(-1, 2)$,
als normaalvector: \omterm{def:g10:algebra:equation}{vergelijking} $-x + 2y + c = 0$; door $P(1,3)$:
$-1 + 6 + c = 0$, dus $c = -5$, en de rechte is $x - 2y + 5 = 0$.

Snijpunt met $d$: uit $2x + y = 7$ volgt $y = 7 - 2x$; invullen geeft
$x - 2(7 - 2x) + 5 = 5x - 9 = 0$, dus $x = \frac95$ en $y = \frac{17}5$:
$H\left(\frac95, \frac{17}5\right)$. Dan is
$\vect{PH}\left(\frac45, \frac25\right)$ en
\[
PH = \sqrt{\frac{16}{25} + \frac{4}{25}} = \frac{\sqrt{20}}{5}
= \frac{2\sqrt5}{5} \approx 0.89 .
\]
\end{solution}

\begin{solution}{exo:g11:scal:10}
Optellen van
\[
\norm{\vec u + \vec v}^2 = \norm{\vec u}^2 + 2\vec u\cdot\vec v +
\norm{\vec v}^2
\quad\text{en}\quad
\norm{\vec u - \vec v}^2 = \norm{\vec u}^2 - 2\vec u\cdot\vec v +
\norm{\vec v}^2
\]
laat de gemengde termen wegvallen en geeft
$\norm{\vec u + \vec v}^2 + \norm{\vec u - \vec v}^2 = 2\norm{\vec u}^2 +
2\norm{\vec v}^2$. In een parallellogram met zijden $\vec u$ en $\vec v$
zijn de diagonalen $\vec u + \vec v$ en $\vec u - \vec v$: de som van de
kwadraten van de twee diagonalen is gelijk aan de som van de kwadraten van
de vier zijden.
\end{solution}

\begin{solution}{exo:g11:scal:11}
\emph{1.} Omdat $I$ het \omterm{prop:g10:coordgeom:midpoint}{midden} is, geldt $\vect{IB} = -\vect{IA}$ en
$IA = 2$. Dan is
\[
\vect{MA}\cdot\vect{MB}
= (\vect{MI} + \vect{IA})\cdot(\vect{MI} - \vect{IA})
= \norm{\vect{MI}}^2 - \norm{\vect{IA}}^2
= MI^2 - 4 .
\]

\emph{2.} De voorwaarde $\vect{MA}\cdot\vect{MB} = 12$ wordt $MI^2 = 16$,
dus $MI = 4$: de verzameling is de cirkel met middelpunt $I$ en straal $4$.
\end{solution}

\begin{solution}{pb:g11:scal:1}
\textbf{1.} $\vec u \cdot \vec v = -3 + 8 = 5$; $\norm{\vec u} = 5$,
$\norm{\vec v} = \sqrt5$; $\cos\theta = \frac{5}{5\sqrt5} =
\frac{1}{\sqrt5}$: $\theta \approx 63.4^\circ$.

\textbf{2.} $6 - 4k = 0$: $k = \frac32$.

\textbf{3.} $W = 50 \times 10 \times \cos 60^\circ = 250$ joule: alleen de
component langs de beweging levert arbeid.

\textbf{4.} $c^2 = 25 + 49 - 2 \times 35 \times \frac12 = 39$:
$c = \sqrt{39}$.

\textbf{5.} $\cos\gamma = \frac{16 + 36 - 64}{2 \times 4 \times 6} =
-\frac14$: negatief, dus is de hoek stomp: $\gamma \approx 104.5^\circ$.

\textbf{6.} \omterm{def:g10:coordgeom:system}{Coördinaten}: $\vec u \cdot \vec v = \cos a \cos b +
\sin a \sin b$. Meetkunde: twee \omterm{def:g11:vect:vector}{vectoren} van lengte $1$, dus
$\vec u \cdot \vec v = 1 \times 1 \times \cos(a - b)$. De twee berekeningen
van hetzelfde getal gelijkstellen geeft
$\cos(a - b) = \cos a \cos b + \sin a \sin b$.

\textbf{7.} $\cos(a + b) = \cos(a - (-b)) = \cos a \cos(-b) +
\sin a \sin(-b) = \cos a \cos b - \sin a \sin b$.

\textbf{8.} $\sin(a + b) = \cos\left(\frac\pi2 - a - b\right)
= \cos\left(\left(\frac\pi2 - a\right) - b\right)
= \cos\left(\frac\pi2 - a\right)\cos b +
\sin\left(\frac\pi2 - a\right)\sin b = \sin a \cos b + \cos a \sin b$.

\textbf{9.} $\cos 15^\circ = \cos 45^\circ \cos 30^\circ +
\sin 45^\circ \sin 30^\circ = \frac{\sqrt2}{2} \cdot \frac{\sqrt3}{2} +
\frac{\sqrt2}{2} \cdot \frac12 = \frac{\sqrt6 + \sqrt2}{4} \approx 0.966$.
En $\sin 75^\circ = \cos 15^\circ$: dezelfde waarde.

\textbf{10.} $b = a$ in de vragen 7 en 8:
$\cos 2a = \cos^2 a - \sin^2 a = 2\cos^2 a - 1 = 1 - 2\sin^2 a$ (met
$\cos^2 + \sin^2 = 1$), en $\sin 2a = 2 \sin a \cos a$.

\textbf{11.} $\sin^2 15^\circ = \frac{1 - \cos 30^\circ}{2} =
\frac{2 - \sqrt3}{4}$, dus $\sin 15^\circ = \frac{\sqrt{2 - \sqrt3}}{2}$.
Het kwadraat van $\frac{\sqrt6 - \sqrt2}{4}$ is
$\frac{6 - 2\sqrt{12} + 2}{16} = \frac{8 - 4\sqrt3}{16} =
\frac{2 - \sqrt3}{4}$: hetzelfde getal in twee kostuums.

\textbf{12.} $\vect{AP}$ valt uiteen in een deel langs $d$ (onzichtbaar voor
$\vec n$) en een deel loodrecht op $d$, waarvan de lengte de gezochte
afstand is; het \omterm{def:g11:scal:dot}{scalair product} met de eenheidsnormaal
$\frac{\vec n}{\norm{\vec n}}$ doodt het eerste deel en meet het tweede. Met
$A$ op $d$: $\vect{AP} \cdot \vec n = a x_0 + b y_0 - (a x_A + b y_A) =
a x_0 + b y_0 + c$, waaruit de formule volgt. Voorbeeld:
$\frac{\abs{15 + 24 - 12}}{5} = \frac{27}{5} = 5.4$.

\textbf{13.} Straal $= 5.4$: $(x - 5)^2 + (y - 6)^2 = 29.16$.

\textbf{14.} Afstand $= \frac{\abs{-12}}{5} = 2.4$. De voet: de loodlijn
door de oorsprong is $(3t, 4t)$; op de rechte geldt $9t + 16t = 12$, dus
$t = \frac{12}{25}$: de voet is $\left(\frac{36}{25}, \frac{48}{25}\right)$,
op afstand $5t = \frac{60}{25} = 2.4$. De formule en de eerlijke berekening
stemmen overeen.

\textbf{15.} Basis $AB = 6$, hoogte $=$ de afstand van $C$ tot $(AB)$ (de as
$y = 0$): $5$; oppervlakte $15$. De rechte $(AC)$ is $5x - 2y = 0$, en
$\operatorname{dist}(B, (AC)) = \frac{\abs{30}}{\sqrt{29}}$; dan is
$\frac12 \times AC \times \frac{30}{\sqrt{29}} = \frac12 \times \sqrt{29}
\times \frac{30}{\sqrt{29}} = 15$: dezelfde oppervlakte, welke hoogtelijn je
ook neemt.

\textbf{16.} $\operatorname{dist} = \frac{\abs{2 + 1 - 6}}{\sqrt2} =
\frac{3}{\sqrt2} \approx 2.12 > 2$: de rechte loopt vrij langs de cirkel ---
zonder gemeenschappelijk punt.

\textbf{17.} Schuif $I$ ertussen:
$\vect{MA} \cdot \vect{MB} = (\vect{MI} + \vect{IA}) \cdot
(\vect{MI} + \vect{IB}) = MI^2 + \vect{MI} \cdot (\vect{IA} + \vect{IB}) +
\vect{IA} \cdot \vect{IB}$. De middelste term sterft
($\vect{IA} + \vect{IB} = \vec 0$), en
$\vect{IA} \cdot \vect{IB} = -IA^2 = -\frac{AB^2}{4}$ (tegengestelde
\omterm{def:g11:vect:vector}{vectoren} van lengte $\frac{AB}2$).

\textbf{18.} $\vect{MA} \cdot \vect{MB} = 0 \iff MI^2 = \frac{AB^2}{4}
\iff MI = \frac{AB}{2}$: de cirkel met middelpunt $I$ en straal
$\frac{AB}2$ --- de cirkel met middellijn $[AB]$.

\textbf{19.} Werk elk paar vanuit $H$ uit ($\vect{BC} = \vect{HC} -
\vect{HB}$, enzovoort):
\[
\vect{HA} \cdot \vect{HC} - \vect{HA} \cdot \vect{HB}
+ \vect{HB} \cdot \vect{HA} - \vect{HB} \cdot \vect{HC}
+ \vect{HC} \cdot \vect{HB} - \vect{HC} \cdot \vect{HA} = 0 :
\]
alles valt paarsgewijs weg --- de identiteit geldt voor \emph{elke} vier
punten. Voor onze $H$ verdwijnen de eerste twee scalaire producten volgens
de veronderstelling, zodat $\vect{HC} \cdot \vect{AB} = 0$: $H$ ligt op de
hoogtelijn uit $C$. Drie hoogtelijnen, één punt.

\textbf{20.} De somformules kwamen uit één \omterm{def:g11:scal:dot}{scalair product}, berekend met het
gezicht van de \emph{\omterm{def:g10:coordgeom:system}{coördinaten}} en dat van de \emph{\omterm{def:g11:scal:dot}{normen} en de hoek}; de
afstandsformule uit het gezicht van de \emph{projectie}; het \omterm{pb:g11:scal:1}{hoogtepunt} uit
het gezicht van de \emph{bilineariteit} (Chasles en \omterm{def:g10:algebra:expand}{uitwerken}). Bereken
dezelfde grootheid twee keer, stel gelijk, en er valt een stelling uit ---
de hele stiel van het \omterm{def:g11:scal:dot}{scalair product}.
\end{solution}
